% Source: https://artofproblemsolving.com/community/c3h1984367_area_of_shaded
% Stable identity: aops-topic-1984367-q01
Put $A=(0,0)$, $B=(12,0)$, $C=(12,12)$, $D=(0,12)$. The arcs centred at $A$ and $B$ meet where $x=6$ and $y=\sqrt{144-36}=6\sqrt3$. By symmetry the region is bounded by four arcs whose endpoints are
\[
(6,6\sqrt3),\quad (6,12-6\sqrt3),\quad (6\sqrt3,6),\quad (12-6\sqrt3,6),
\]
the vertices of a square (dashed) whose diagonals have length $6\sqrt3-(12-6\sqrt3)=12(\sqrt3-1)$.

\begin{asy}
usepackage("amsmath");
size(5.5cm);
defaultpen(fontsize(10pt));
pair A=(0,0),B=(12,0),C=(12,12),D=(0,12);
path q1=arc(A,12,0,90),q2=arc(B,12,90,180),q3=arc(C,12,180,270),q4=arc(D,12,270,360);
real s=6*sqrt(3);
pair T=(6,s),Bo=(6,12-s),Ri=(s,6),Le=(12-s,6);
path reg=arc(D,12,-90+30,-90+60)--arc(A,12,30,60)--arc(B,12,120,150)--arc(C,12,210,240)--cycle;
fill(reg,gray(0.75));
draw(A--B--C--D--cycle);
draw(q1,gray(0.4));draw(q2,gray(0.4));draw(q3,gray(0.4));draw(q4,gray(0.4));
draw(Bo--Ri--T--Le--cycle,dashed);
dot(T);dot(Bo);dot(Ri);dot(Le);
label("$A$",A,SW);label("$B$",B,SE);label("$C$",C,NE);label("$D$",D,NW);
label("$12$",(A+B)/2,S);
label("\mbox{}",(min(currentpicture,user=true)+max(currentpicture,user=true))/2,black); // keeps the text after the figure black when LaTeX typesets the labels inline
\end{asy}

\textbf{Inner square.} Its area is half the square of a diagonal: $\dfrac{144(\sqrt3-1)^2}{2}=72(4-2\sqrt3)=288-144\sqrt3$.

\textbf{Four circular segments.} The side from $(12-6\sqrt3,6)$ to $(6,6\sqrt3)$ is a chord of the circle centred at $B$; both endpoints make angles of $150^\circ$ and $120^\circ$ with the positive $x$-axis as seen from $B$, so the chord subtends $30^\circ$ at $B$. A segment of radius $12$ and angle $\dfrac{\pi}{6}$ has area
\[
\dfrac{12^2}{2}\left(\dfrac{\pi}{6}-\sin\dfrac{\pi}{6}\right)=12\pi-36 .
\]
The total area is
\[
288-144\sqrt3+4(12\pi-36)=\boxed{144+48\pi-144\sqrt3}\approx45.4 .
\]
